% =========================================================
% Aporte individual: Alcivar Velez Anderson Adonis
% Seccion: Gestion del Proyecto, Riesgos y Declaracion de uso de IA
% Entrega: TA-PFC-E1 -- DentiTech (Clinica Odontologica Medicus)
% Citas usadas (ver bibliografia_ieee.bib / referencias.bib):
%   pmbok2021, iso29148, pohl2025, wiegers2023, iso31000, lopdp2021
% =========================================================

\section{Gesti{\'o}n del Proyecto}

El Proyecto Fin de Curso (PFC) \textit{DentiTech} se desarrolla para la Cl{\'i}nica Odontol{\'o}gica \textit{Medicus} (Cooperativa Jet, Quevedo) durante el per{\'i}odo acad{\'e}mico SPA 2026-2027. Su gesti{\'o}n se basa en los principios de entrega de valor, liderazgo compartido y adaptabilidad descritos en la gu{\'i}a PMBOK \cite{pmbok2021}, y en los procesos de ingenier{\'i}a de requisitos de la norma ISO/IEC/IEEE 29148 \cite{iso29148}. Esta secci{\'o}n define qui{\'e}n hace qu{\'e}, cu{\'a}ndo se produce cada entregable, c{\'o}mo se coordina el equipo con la cl{\'i}nica y qu{\'e} riesgos podr{\'i}an afectar el proyecto.

\subsection{Asignaci{\'o}n de Roles del Equipo}

El equipo est{\'a} formado por cuatro estudiantes de la carrera de Ingenier{\'i}a en Software de la UTEQ. Cada integrante tiene una responsabilidad principal y una carpeta propia en \texttt{PFC/E1/aportes/}, de modo que su contribuci{\'o}n individual queda registrada en el historial de GitHub. La Tabla~\ref{tab:and_roles} resume la asignaci{\'o}n.

\begin{table}[htbp]
\centering
\caption{Roles y responsabilidades del equipo DentiTech}
\label{tab:and_roles}
\begin{tabularx}{\textwidth}{>{\raggedright\arraybackslash}p{3.6cm} >{\raggedright\arraybackslash}p{3.6cm} >{\raggedright\arraybackslash}X}
\toprule
\textbf{Integrante} & \textbf{Rol} & \textbf{Responsabilidades en E1} \\
\midrule
Alc{\'i}var V{\'e}lez Anderson Adonis & Gesti{\'o}n del Proyecto, Riesgos y Referencias & Roles, cronograma (semanas 4 a 17), plan de comunicaci{\'o}n con la cl{\'i}nica, matriz de riesgos, declaraci{\'o}n de uso de IA y bibliograf{\'i}a IEEE. \\
\addlinespace
Contreras Ch{\'a}vez Kevin Germ{\'a}n & L{\'i}der de Repositorio, Visi{\'o}n, Alcance e Integraci{\'o}n & Ficha del proyecto, visi{\'o}n, objetivos de negocio medibles, alcance y restricciones, carta de aceptaci{\'o}n, manifiesto SHA-256 e integraci{\'o}n del informe en LaTeX. \\
\addlinespace
Fuertes Calvo Landy Lorena & Partes Interesadas y Contexto del Sistema & Identificaci{\'o}n de partes interesadas por rol, matriz poder-inter{\'e}s, diagrama de contexto y an{\'a}lisis de las ocho perspectivas. \\
\addlinespace
Mor{\'a}n Pilaguano Frixon Fernando & Elicitaci{\'o}n, Instrumentos y {\'E}tica LOPDP & Plan de elicitaci{\'o}n, instrumentos de levantamiento, protocolo de anonimizaci{\'o}n y consentimiento informado. \\
\bottomrule
\end{tabularx}
\end{table}

Adem{\'a}s de su rol principal, todos los integrantes participan en las sesiones de elicitaci{\'o}n en la cl{\'i}nica y revisan el informe integrado antes de cada entrega. Las decisiones que afectan al alcance se toman por consenso del equipo y quedan registradas en el repositorio.

